\medskip

\begin{minipage}{6.7cm}On considère le cube ABCDEFGH d'arête 1 représenté ci-contre.

On note K le milieu du segment [HG].

On se place dans le repère orthonormé $\left(\text{A}~;~\vect{\text{AD}},~\vect{\text{AB}}, \vect{\text{AE}}\right)$.

\smallskip

\begin{enumerate}
\item Justifier que les points C, F et K définissent un plan.
\item 
	\begin{enumerate}
		\item Donner, sans justifier, les longueurs KG, GF et GC.
		\item Calculer l'aire du triangle FGC.
		\item Calculer le volume du tétraèdre FGCK.
	\end{enumerate}
\end{enumerate}	
\end{minipage}\hfill
\begin{minipage}{6.7cm}
\psset{unit=0.75cm,arrowsize=2pt 3}
\begin{pspicture}(-2.8,-2.4)(5,5.3)
\psline[linestyle=dashed](3,0)(0,0)(0,3)
\psline[linestyle=dashed](-1.3,-1.3)
\psline[linewidth=1.25pt](3,0)(5,0)
\psline[linewidth=1.25pt](0,3)(0,5.3)
\psline[linewidth=1.25pt](-1.3,-1.3)(-2.8,-2.8)
\psframe[linewidth=1.25pt](-1.3,-1.3)(1.7,1.7)
\psline[linewidth=1.25pt](1.7,-1.3)(3,0)(3,3)(1.7,1.7)
\psline[linewidth=1.25pt](3,3)(0,3)(-1.3,1.7)
\uput[ur](0,0){A} \uput[dr](3,0){B} \uput[dr](1.7,-1.3){C} \uput[dr](-1.3,-1.3){D}
\uput[ul](0,3){E} \uput[ur](3,3){F} \uput[ul](1.7,1.7){G} \uput[ul](-1.3,1.7){H}
\uput[ul](-2.8,-2.8){$x$} \uput[u](5,0){$y$} \uput[r](0,5.3){$z$}
\psdots(0,0)(3,0)(1.7,-1.3)(-1.3,-1.3)(0,3)(3,3)(1.7,1.7)(-1.3,1.7)(0.2,1.7)
\uput[u](0.2,1.7){K}
\end{pspicture}
\end{minipage}

\smallskip
	
On rappelle que le volume $V$ d'un tétraèdre est donné par:
\[V = \dfrac13\mathcal{B} \times h,\]
où $\mathcal{B}$ est l'aire d'une base et $h$ la hauteur correspondante.

\begin{enumerate}[start=3]
\item 
	\begin{enumerate}
		\item On note $\vect{n}$ le vecteur de coordonnées $\begin{pmatrix}1\\2\\1\end{pmatrix}$.
		
Démontrer que $\vect{n}$ est normal au plan (CFK).
		\item  En déduire qu'une équation cartésienne du plan (CFK) est:
\[x + 2y + z - 3 = 0.\]
	\end{enumerate}
\item On note $\Delta$ la droite passant par le point G et orthogonale au plan (CFK). 

Démontrer qu'une représentation paramétrique de la droite $\Delta$ est:
\[\left\{\begin{array}{l c l}
x&=&1 + t\\
y&=&1 + 2t\\
z&=&1 + t
\end{array}\right. \quad  (t \in \R))\]
\item Soit L le point d'intersection entre la droite $\Delta$ et le plan (CFK).
	\begin{enumerate}
		\item Déterminer les coordonnées du point L.
		\item En déduire que LG $= \dfrac{\sqrt 6}{6}$.
	\end{enumerate}
\item En utilisant la question 2., déterminer la valeur exacte de l'aire du triangle CFK.
\end{enumerate}

\bigskip
